% Generated draft: requires mathematical and source review.
\documentclass[11pt]{article}
\usepackage[margin=25mm]{geometry}
\usepackage{fontspec}
\setmainfont{Latin Modern Roman}
\usepackage{amsmath,amssymb,mathtools,unicode-math}
\setmathfont{Latin Modern Math}
\usepackage{enumitem,needspace,xurl,hyperref}
\setlength{\parindent}{0pt}
\setlength{\parskip}{5pt}
\setlength{\emergencystretch}{3em}
\newcommand{\N}{\mathbb{N}}
\newcommand{\Z}{\mathbb{Z}}
\newcommand{\Q}{\mathbb{Q}}
\newcommand{\R}{\mathbb{R}}
\newcommand{\C}{\mathbb{C}}
\newcommand{\F}{\mathbb{F}}
\newcommand{\A}{\mathbb{A}}
\newcommand{\PP}{\mathbb{P}}
\newcommand{\E}{\mathbb{E}}
\newcommand{\eps}{\varepsilon}
\newcommand{\dd}{\,\mathrm{d}}
\begin{document}
\section*{UnsolvedMath ID: 30003367}
\textbf{Source problem code:} catalogue entry
\subsection{Definitions and mathematical statement}
A conventional framework for interpreting the supplied questions is the following; the source record does not specify that this is the intended full set of hypotheses. Let \(M\) be a closed smooth manifold, let \(\operatorname{Diff}^{\infty}(M)\) denote its group of smooth diffeomorphisms, and fix an auxiliary Riemannian norm on \(TM\). A diffeomorphism \(f\in\operatorname{Diff}^{\infty}(M)\) is Anosov if there exist a continuous \(Df\)-invariant splitting \(TM=E^s\oplus E^u\), a constant \(C\geq 1\), and \(0<\lambda<1\) such that, for every \(x\in M\), every integer \(n\geq0\), and all \(v^s\in E_x^s\), \(v^u\in E_x^u\),
\[
\|D(f^n)_xv^s\|\leq C\lambda^n\|v^s\|,
\qquad
\|D(f^{-n})_xv^u\|\leq C\lambda^n\|v^u\|.
\]
Write \(\mathcal A^{\infty}(M)\) for the set of such diffeomorphisms. Two elements \(f,g\in\mathcal A^{\infty}(M)\) are smoothly conjugate if there is \(h\in\operatorname{Diff}^{\infty}(M)\) with \(g=h\circ f\circ h^{-1}\). Define the smooth-conjugacy class and the corresponding set of equivalence classes by
\[
\mathcal O^{\infty}(f)=\{h\circ f\circ h^{-1}:h\in\operatorname{Diff}^{\infty}(M)\},
\qquad
\mathcal M_A^{\infty}(M)=\mathcal A^{\infty}(M)/{\sim},
\]
where \(\sim\) is smooth conjugacy.

For a metric space \((X,d)\) and \(S\subseteq X\), define, for \(s>0\),
\[
\mathcal H_d^s(S)=\lim_{\delta\downarrow0}\inf\left\{
\sum_{i=1}^{\infty}(\operatorname{diam}_d U_i)^s:
S\subseteq\bigcup_{i=1}^{\infty}U_i,\quad
\operatorname{diam}_d U_i\leq\delta
\right\},
\qquad
\dim_{H,d}S=\inf\{s>0:\mathcal H_d^s(S)=0\},
\]
with \(\inf\varnothing=+\infty\).

\textbf{Questions recorded by the source.} Can a smooth-conjugacy class of Anosov diffeomorphisms have Hausdorff dimension greater than \(2\)? Can Ebin's infinite-dimensional analysis of Riemannian metrics modulo diffeomorphisms be adapted to Anosov diffeomorphisms and their moduli spaces?

\textbf{Unresolved formulation.} Under the provisional framework above, the first question would ask whether \(\dim_{H,d}\mathcal O^{\infty}(f)>2\) can occur for an admissible \(M\) and \(f\), but the record supplies neither the metric \(d\) nor the admissible manifolds, their dimensions, or the intended differentiability classes. Hausdorff dimension depends on the metric; a topology alone does not determine it. Thus this expression does not yet specify a unique mathematical problem. The second question is a research-program question: the record does not identify the particular Ebin results or analytical structures to be reproduced, the topology or geometry of the moduli space, or the criterion for a successful adaptation. The definitions above provide terminology and do not resolve these omissions. The original statement contains only the second question, whereas the supplied rewritten statements also contain the Hausdorff-dimension question.
\subsection{Short English statement}
The supplied record asks whether a smooth-conjugacy class of Anosov diffeomorphisms can have Hausdorff dimension greater than \(2\), and whether Ebin's methods for Riemannian metrics modulo diffeomorphisms extend to this setting. A precise statement requires the intended metric, underlying manifolds, regularity, and analytical goals. The record labels the problem open; that label is supplied curation and has not been independently verified.
\subsection{Sources}
\begin{itemize}
\item Supplied local UnsolvedMath registry record, permanent ID \(30003367\), title \emph{Ebin Analysis for Anosov Diffeomorphism Moduli}, exact source problem code \texttt{catalogue entry}, in the Oberwolfach Reports Open Problems collection. The record gives proposed year \(2017\) and names no proposer.
\item Supplied source citation: \emph{Mini-Workshop: Spaces and Moduli Spaces of Riemannian Metrics (2018)}, Oberwolfach Reports, pp.\ 133--166, DOI \texttt{10.4171/owr/2017/3}. \url{https://doi.org/10.4171/owr/2017/3}. The citation's parenthetical year \(2018\) differs from the record's proposed year and its identification as Report 3/2017; this discrepancy remains unresolved.
\item Supplied literature source: \emph{Oberwolfach Report 3/2017}. \url{https://ems.press/content/serial-article-files/46667}.
\item Supplied literature source: \emph{Smooth conjugacy of Anosov diffeomorphisms}. \url{https://arxiv.org/abs/0804.3901}.
\item Literature assessment dated exactly \texttt{2026-08-21}: the supplied curation labels the problem open and reports that no general Ebin-style infinite-dimensional analysis for the Anosov setting is known, while existing work parametrizes special smooth-conjugacy classes and establishes rigidity only in restricted families. This is supplied dated curation, not an independently verified literature assessment. The background explicitly requests independent verification of time-sensitive status and cited claims.
\item The registry assigns \texttt{statement\_status: corrected\_verified} and reports that the official passage was checked and rewritten while preserving its mathematical content. These are supplied provenance claims; no independent source inspection is asserted here. Ebin is named in the problem, but no specific bibliographic reference to his analysis is supplied.
\end{itemize}
\end{document}
